\section{The Unit Store}
\label{sec:sec03}
\label{sec:unit-store}
%==============================================================================
\newtheorem*{condition}{Condition}

Section~\ref{sec:ledger} fixed the unit universe $\mathcal{U}$ and the balance function $w_t : \mathcal{U} \to \mathbb{R}$ but left $\mathcal{U}$ abstract. The \textbf{Unit Store} makes it concrete: it decides which units exist, records their identity and static terms, binds each to a smart contract, and validates every unit at registration. The executor, the settlement projection, and the valuation function all read $\mathcal{U}$ from here.

\subsection{Units}

A unit is any financial instrument that can appear as a wallet balance: cash currencies, listed equities and bonds (by ISIN), listed derivatives (by exchange contract specification), OTC derivatives (each bilateral trade), and structured products (each issuance). A unit has \emph{identity} (how it is distinguished from every other unit), \emph{type}, static \emph{terms}, and a bound \emph{smart contract}. Its evolving state---lifecycle stage and product-specific state---is not stored on the unit entry; it lives in the three-home state model (Section~\ref{sec:sec04}).

\subsection{Unit Identity: OTC versus Listed}
\label{sec:sec03-identity}
\label{sec:unit-identity}

Two positions are in the same unit exactly when they are fungible.

\begin{principle}[Unit Identity]
\label{prin:unit-identity}
Two positions are in the same unit if and only if they are fungible: a holder is economically indifferent between them. Fungible positions share one unit entry and net against each other; non-fungible positions are distinct units with distinct entries.
\end{principle}

Fungibility fixes a different identity key per instrument class.

\begin{center}
\small
\begin{tabular}{lll}
\toprule
\textbf{Instrument class} & \textbf{Identity key} & \textbf{Fungibility} \\
\midrule
Cash & Currency code & All balances of a currency net \\
Listed equity & ISIN & All shares of an ISIN net \\
Listed derivative & Contract spec & All positions in a series net \\
 & (exchange, underlier, type, strike, expiry) & (via CCP novation) \\
OTC derivative & Full CDM \texttt{Trade} & Non-fungible: counterparty and \\
 & (includes counterparty, \texttt{Collateral}) & CSA make each trade unique \\
Bond & ISIN & All holdings of an ISIN net \\
Structured note & ISIN + product terms & Same issuance is fungible \\
Tokenised equity & Contract address + chain ID & All tokens on a contract net \\
\bottomrule
\end{tabular}
\end{center}

For OTC instruments the identity is the CDM \texttt{Trade} object \emph{including} its \texttt{Collateral} field: two trades with identical payoffs under different CSAs (USD versus EUR) are different units, with different discount curves, margin flows, and settlement currencies. For listed instruments the CCP novates itself in as sole counterparty, so all positions in one contract specification are fungible regardless of original counterparty. A tokenised security is a unit distinct from its underlying---NVDA and tokenised-NVDA carry the same exposure under different settlement mechanisms; the double-counting resolution (the custodian nets to zero economic exposure) is in Appendix~\ref{app:cdm-walkthrough}.

\subsection{Three-Tier Architecture}

The store has three tiers, each answering one question.

\paragraph{Tier 1 --- Reference Data. \emph{What instruments exist in the market?}}
Instrument master data---ISIN, exchange, contract specification, issuer, coupon schedule---read from exchanges, vendors, and CSD records. The Ledger consumes reference data; it does not author it. CDM: \texttt{TransferableProduct} for securities, \texttt{NonTransferableProduct} for derivatives and structured products.

\paragraph{Tier 2 --- Product Registry. \emph{What lifecycle logic governs this kind of instrument?}}
One smart-contract template per product \emph{type}, not per contract: ``European equity index option'' maps to one \texttt{OptionPayout} template covering every such option regardless of underlier, strike, or expiry. CDM \texttt{ProductQualification} classifies \texttt{EconomicTerms} into a product type; each type maps to a template (Section~\ref{sec:contracts}). An entry is created automatically the first time a Tier-3 registration meets a new product type, and is thereafter immutable.

\paragraph{Tier 3 --- Unit Registry. \emph{Which specific units can appear as balances?}}
The elements of $\mathcal{U}$. For listed instruments the unit is the contract specification; for OTC instruments it is the full CDM \texttt{Trade}, \texttt{Collateral} included. Each entry carries identity, type, the static terms drawn from \texttt{EconomicTerms}, the bindings to its Tier-2 template and smart contract, and provenance:

\begin{lstlisting}
UnitEntry:
    unit_id        : UniqueId        -- deterministic, from the CDM object
    unit_type      : UnitType        -- CASH | EQUITY | LISTED_DERIV
                                     --   | OTC_DERIV | BOND | STRUCTURED
    -- Identity (varies by type)
    isin           : Option<ISIN>          -- securities, bonds, notes
    contract_spec  : Option<ContractSpec>  -- listed derivatives
    cdm_trade_ref  : Option<TradeRef>      -- OTC (includes Collateral)
    -- Bindings
    product_ref    : ProductRegistryRef    -- Tier 2 template
    smart_contract : SmartContractRef      -- lifecycle logic instance
    -- Static terms (from CDM EconomicTerms)
    currency       : CurrencyCode
    multiplier     : Option<Decimal>
    expiry         : Option<Date>
    -- Provenance
    created_by     : TransactionId         -- registering transaction
    created_at     : Timestamp
\end{lstlisting}

The \texttt{unit\_id} is derived deterministically and injectively from the CDM object: a hash of the contract-specification fields for listed instruments, the UTI or metadata key for OTC. Distinct CDM objects always yield distinct unit IDs.

\paragraph{Where evolving state lives.}
The unit entry holds identity and static terms only. The lifecycle stage and the product-specific state---which coupons a bond has paid, the accumulated cost of a futures position---are not entry fields. A unit carries immutable, versioned terms in \emph{ProductTerms}$[u]$ and a shared lifecycle status in \emph{UnitStatus}$[u]$; per-position state lives in \emph{PositionState}$[w,u]$. These are the three-home state model (Section~\ref{sec:sec04}). The lifecycle stage is read from \emph{UnitStatus}.

\begin{remark}[UnitStatus is a projection, not a mutable store]
\label{rem:sec03-unitstatus}
UnitStatus holds one value per unit, shared---read identically by every holder. Its value changes over time, but UnitStatus is not a separate source of truth: the immutable event log is. Every change is caused by a logged event; the stored value is overwritten in place only as a read cache. Replaying the events up to any point rebuilds the exact value that held then, so nothing is lost by overwriting and there is no other way to change the value. A value exists from registration onward.
\end{remark}

\subsection{Registration}
\label{sec:sec03-registration}
\label{sec:unit-registration}

Units enter the registry through three channels.

\begin{enumerate}
\item \textbf{Cash (pre-registered).} Currency units are permanent entries created at system initialisation: no expiry, no lifecycle events, stage \texttt{ACTIVE} permanently, smart contract the identity function.
\item \textbf{Listed instruments (reference-data feed).} On a new listing or initialisation against an existing market, the feed supplies the contract specification; the store creates the Tier-1, Tier-2 (if the product type is new), and Tier-3 entries. This is administrative, not a trade: the unit enters $\mathcal{U}$ upon listing, before any position is taken.
\item \textbf{OTC instruments (trade execution).} The executed CDM \texttt{Trade} \emph{is} the unit definition; the store creates the Tier-3 entry from it. The unit ID derives from the Trade's metadata key (counterparty and \texttt{Collateral} included), so identical payoffs under different CSAs receive different IDs. The unit enters $\mathcal{U}$ at execution.
\end{enumerate}

Registration is atomic: a unit is fully validated and added, or rejected with no partial state. Units are never deleted; a unit at end of life reaches a terminal stage in \emph{UnitStatus} (matured, terminated, settled) and the entry remains for audit and time travel, while the executor rejects new moves against a non-\texttt{ACTIVE} unit except to settle existing obligations.

Two disciplines bind registration. Registration writes the first value of every per-unit map, and it happens once.

\begin{condition}[C7: ProductTerms registration-total]
\label{cond:C7}
\sloppy
\emph{ProductTerms} is registration-total: its first version is written at registration, so \texttt{product\_terms(u)} is defined for every registered $u$. (\emph{UnitStatus} is registration-total by the same write, C5, established in Section~\ref{sec:sec04}.)
\end{condition}

\begin{condition}[C10: no re-registration]
\label{cond:C10}
Re-registering an existing unit id is a hard error, never a silent reset of its state.
\end{condition}

Registration writes \emph{ProductTerms} and \emph{UnitStatus} together, so ``registered in ProductTerms $\iff$ registered in UnitStatus'' holds by construction; C10 makes the write once-only, so a re-registration cannot overwrite live state behind the log.

The registering function writes the \emph{singleton} first version of \emph{ProductTerms} and the default \emph{UnitStatus} together, returning the updated ledger, or rejecting an already-registered id with a \texttt{Left} (C10). \emph{ProductTerms} is a \texttt{NonEmpty} list of versions, so a registered-but-versionless unit is unrepresentable and the in-force-version accessor is total---it needs no \texttt{Maybe} for emptiness (C7). The two-map write is a single operation, so the coupling invariant ``registered in \emph{ProductTerms} $\iff$ registered in \emph{UnitStatus}'' cannot be falsified by any caller; \texttt{register} is exported, but \texttt{ProductTerms} and \texttt{Ledger} are exported abstract (no field setter), so the singleton is the only origin of a unit's terms and there is no second door.

\begin{lstlisting}[language=Haskell]
-- ProductTerms abstract: NonEmpty makes "registered but versionless"
-- unrepresentable; currentTerms is then total (no Maybe).      -- C7
newtype ProductTerms = ProductTerms (NonEmpty TermsVersion)

data LedgerError = ReRegistration UnitId | UnknownUnit UnitId    -- C10

-- Registration is ONE move-less Transaction whose txIntroduce carries
-- the unit's first terms version and its default status.
registerTx :: UnitId -> TermsVersion -> UnitStatus -> Transaction
registerTx u tv us = Transaction
  { txUnit = u, txMoves = [], txRows = Map.empty, txStatus = []
  , txIntroduce = Just (tv, us), txAppend = Nothing }

-- register is not a separate door: it is applyTx of that transaction, so
-- registration, trade, settlement, and amendment share one apply. The
-- introduce branch of applyTx writes the singleton first ProductTerms
-- version (C7) and the default UnitStatus (C5) together, and rejects an
-- already-registered id with Left (ReRegistration u) (C10). [Section 4]
register :: UnitId -> TermsVersion -> UnitStatus
         -> Ledger -> Either LedgerError Ledger
register u tv us = applyTx (registerTx u tv us)
\end{lstlisting}

A first registration succeeds; the same id a second time is rejected, leaving the prior ledger and its live state untouched---a hard error, never a silent reset:

\begin{lstlisting}[language=Haskell]
register uES tvES  defaultStatus emptyLedger  ==  Right l0
register uES tvDup defaultStatus l0           ==  Left (ReRegistration uES)
\end{lstlisting}

Conservation does not arise at registration: a registration is a move-less transaction, so its edge-sum is \texttt{mempty}---vacuously conserved---and it touches no \emph{PositionState} row. A non-vacuous balance, and the conservation that constrains it, appears only when the first position is taken (Section~\ref{sec:sec04}).

\subsection{Two-Stage Validation}
\label{sec:sec03-validation}
\label{sec:unit-validation}

Validity is enforced at registration and at transaction time.

\noindent\textbf{At registration}, a unit must satisfy:
\begin{enumerate}[nosep]
\item \textbf{Uniqueness}: no existing unit shares its \texttt{unit\_id} (C10).
\item \textbf{Smart-contract binding}: a contract reference exists for its product type.
\item \textbf{CDM qualification}: for CDM-native units, \texttt{ProductQualification} returns a valid classification for the stated \texttt{EconomicTerms}; an unclassifiable product is rejected.
\item \textbf{Term consistency}: static terms cohere---expiry in the future at creation, positive multiplier, valid currency code.
\end{enumerate}

\noindent\textbf{At transaction time}, the executor checks that every unit a move references (i) exists in the store, (ii) is in a stage compatible with the operation (read from \emph{UnitStatus}), and (iii) has a smart contract invokable for the requested event.

Invalid units cannot enter $\mathcal{U}$, and moves cannot reference non-existent or retired units. With the conservation law (Section~\ref{sec:ledger}) and the executor's purity (Section~\ref{sec:lifecycle}), this completes the correctness chain: the ledger operates only on validated, contract-governed units.

\subsection{Guarantees}

The Unit Store provides four guarantees to every downstream consumer.

\begin{enumerate}
\item \textbf{Existence}: every unit a move references exists and is in a valid lifecycle stage.
\item \textbf{Immutability of identity and terms}: a \texttt{unit\_id} always denotes the same instrument with the same static terms. Lifecycle stage and product-specific state evolve (Section~\ref{sec:sec04}); identity and terms do not.
\item \textbf{Smart-contract availability}: every registered unit has a bound contract the executor can invoke for lifecycle events.
\item \textbf{CDM alignment}: every CDM-native unit maps to a CDM object (\texttt{Trade}, \texttt{Non\allowbreak{}Transferable\allowbreak{}Product}, or \texttt{Transferable\allowbreak{}Product}) with a successful qualification.
\end{enumerate}

These make referential integrity (P3, Section~\ref{sec:invariants}) hold by construction: $m.\text{unit} \in \mathcal{U}$ because the executor consults the store before accepting any move.

\subsection{CDM Alignment}

Each tier maps to a CDM concept.

\begin{center}
\small
\begin{tabular}{lll}
\toprule
\textbf{Tier} & \textbf{Content} & \textbf{CDM equivalent} \\
\midrule
1 Reference Data & Instrument master data & \texttt{TransferableProduct} (securities), \\
 & & \texttt{NonTransferableProduct} (derivatives) \\
2 Product Registry & Smart-contract templates & \texttt{ProductQualification} + \texttt{EconomicTerms} \\
3 Unit Registry & Position-bearing units ($\mathcal{U}$) & \texttt{Trade} (OTC), contract spec (listed) \\
\bottomrule
\end{tabular}
\end{center}

CDM has no native catalogue concept: it models individual trades, not registries, and serves both listed and OTC derivatives through one \texttt{NonTransferableProduct} type without a dedicated listed-contract-specification type or a model for exchange reference-data feeds (listings, expiry calendars, lot sizes). These are documentation gaps, not design incompatibilities: the store uses CDM objects as its internal representation, and catalogue management sits alongside CDM as a Ledger-specific concern.
